% Abreviaturas empleadas en el texto principal y en la metodología.
\newacronym{acr-ia}{IA}{inteligencia artificial (\textit{Artificial Intelligence})}
\newacronym{acr-ml}{ML}{aprendizaje automático (\textit{Machine Learning})}
\newacronym{acr-sl}{SL}{aprendizaje supervisado (\textit{Supervised Learning})}
\newacronym{acr-ul}{UL}{aprendizaje no supervisado (\textit{Unsupervised Learning})}
\newacronym{acr-rl}{RL}{aprendizaje por refuerzo (\textit{Reinforcement Learning})}
\newacronym{acr-marl}{MARL}{aprendizaje por refuerzo multi-agente (\textit{Multi-Agent Reinforcement Learning})}
\newacronym{acr-madrl}{MADRL}{aprendizaje por refuerzo profundo multi-agente (\textit{Multi-Agent Deep Reinforcement Learning})}
\newacronym{acr-harl}{HARL}{aprendizaje por refuerzo con agentes heterogéneos (\textit{Heterogeneous-Agent Reinforcement Learning})}
\newacronym{acr-il}{IL}{aprendizaje independiente (\textit{Independent Learning})}
\newacronym{acr-mdp}{MDP}{proceso de decisión de Markov (\textit{Markov Decision Process})}
\newacronym{acr-pomdp}{POMDP}{proceso de decisión de Markov parcialmente observable (\textit{Partially Observable Markov Decision Process})}
\newacronym{acr-decpomdp}{Dec-POMDP}{proceso de decisión de Markov parcialmente observable descentralizado (\textit{Decentralized Partially Observable Markov Decision Process})}
\newacronym{acr-ctde}{CTDE}{entrenamiento centralizado con ejecución descentralizada (\textit{Centralized Training with Decentralized Execution})}
\newacronym{acr-dqn}{DQN}{red neuronal profunda de acción-valor (\textit{Deep Q-Network})}
\newacronym{acr-idqn}{IDQN}{DQN independiente (\textit{Independent Deep Q-Network})}
\newacronym{acr-drl}{DRL}{aprendizaje por refuerzo profundo (\textit{Deep Reinforcement Learning})}
\newacronym{acr-td}{TD}{diferencia temporal (\textit{Temporal-Difference})}
\newacronym{acr-gae}{GAE}{estimación generalizada de la ventaja (\textit{Generalized Advantage Estimation})}
\newacronym{acr-trpo}{TRPO}{optimización de política mediante región de confianza (\textit{Trust Region Policy Optimization})}
\newacronym{acr-ppo}{PPO}{optimización de política próxima (\textit{Proximal Policy Optimization})}
\newacronym{acr-ippo}{IPPO}{optimización de política próxima independiente (\textit{Independent Proximal Policy Optimization})}
\newacronym{acr-mappo}{MAPPO}{optimización de política próxima multi-agente (\textit{Multi-Agent Proximal Policy Optimization})}
\newacronym{acr-happo}{HAPPO}{optimización de política próxima para agentes heterogéneos (\textit{Heterogeneous-Agent Proximal Policy Optimization})}
\newacronym{acr-hatrpo}{HATRPO}{optimización de política mediante región de confianza para agentes heterogéneos (\textit{Heterogeneous-Agent Trust Region Policy Optimization})}
\newacronym{acr-a2c}{A2C}{ventaja del actor-crítico (\textit{Advantage Actor-Critic})}
\newacronym{acr-ddpg}{DDPG}{gradiente determinista de política profunda (\textit{Deep Deterministic Policy Gradient})}
\newacronym{acr-maddpg}{MADDPG}{gradiente determinista de política profunda multi-agente (\textit{Multi-Agent Deep Deterministic Policy Gradient})}
\newacronym{acr-qmix}{QMIX}{red de mezcla de valores de acción (\textit{Q-Mixing Network})}
\newacronym{acr-vdn}{VDN}{redes de descomposición de valor (\textit{Value-Decomposition Networks})}
\newacronym{acr-phc}{PHC}{ascenso de colinas de política (\textit{Policy Hill-Climbing})}
\newacronym{acr-wolf-phc}{WoLF-PHC}{ascenso de colinas de política con principio ganar o aprender rápido (\textit{Win or Learn Fast Policy Hill-Climbing})}
\newacronym{acr-mlp}{MLP}{perceptrón multicapa (\textit{Multi-Layer Perceptron})}
\newacronym{acr-relu}{ReLU}{unidad lineal rectificada (\textit{Rectified Linear Unit})}
\newacronym{acr-sgd}{SGD}{descenso de gradiente estocástico (\textit{Stochastic Gradient Descent})}
\newacronym{acr-api}{API}{interfaz de programación de aplicaciones (\textit{Application Programming Interface})}
\newacronym{acr-xml}{XML}{lenguaje de marcado extensible (\textit{Extensible Markup Language})}
\newacronym{acr-tcp}{TCP}{protocolo de control de transmisión (\textit{Transmission Control Protocol})}
\newacronym{acr-uml}{UML}{lenguaje unificado de modelado (\textit{Unified Modeling Language})}
\newacronym{acr-gui}{GUI}{interfaz gráfica de usuario (\textit{Graphical User Interface})}
\newacronym{acr-osm}{OSM}{OpenStreetMap}
\newacronym{acr-sumo}{SUMO}{Simulación de Movilidad Urbana (\textit{Simulation of Urban MObility})}
\newacronym{acr-dlr}{DLR}{Centro Aeroespacial Alemán (\textit{Deutsches Zentrum für Luft- und Raumfahrt})}
\newacronym{acr-traci}{TraCI}{interfaz de control del tránsito (\textit{Traffic Control Interface})}
\newacronym{acr-taz}{TAZ}{zona de análisis de tránsito (\textit{Traffic Analysis Zone})}
\newacronym{acr-od}{OD}{origen-destino (\textit{Origin-Destination})}
\newacronym{acr-bfs}{BFS}{búsqueda en anchura (\textit{Breadth-First Search})}
\newacronym{acr-dag}{DAG}{grafo acíclico dirigido (\textit{Directed Acyclic Graph})}
\newacronym{acr-e1}{E1}{detector de bucle inductivo puntual de entrada}
\newacronym{acr-e2}{E2}{detector de área continua de carril}
\newacronym{acr-hbefa}{HBEFA}{manual de factores de emisión para transporte por carretera (\textit{Handbook Emission Factors for Road Transport})}
\newacronym{acr-phemlight}{PHEMlight}{modelo microscópico de emisiones de vehículos livianos y pesados (\textit{Passenger car and Heavy duty Emission Model light})}
\newacronym{acr-hpo}{HPO}{optimización de hiperparámetros (\textit{Hyperparameter Optimization})}
\newacronym{acr-tpe}{TPE}{estimador de densidad de Parzen estructurado por árbol (\textit{Tree-structured Parzen Estimator})}
\newacronym{acr-asha}{ASHA}{algoritmo de reducción sucesiva de configuraciones asíncrono (\textit{Asynchronous Successive Halving Algorithm})}
\newacronym{acr-hpc}{HPC}{computación de alto rendimiento (\textit{High-Performance Computing})}
\newacronym{acr-nfs}{NFS}{sistema de archivos de red (\textit{Network File System})}
\newacronym{acr-slurm}{SLURM}{gestor simple de recursos para administración de Linux (\textit{Simple Linux Utility for Resource Management})}
\newacronym{acr-ccar}{CCAR}{Centro de Cómputo de Alto Rendimiento}
\newacronym{acr-ccad}{CCAD}{Centro de Cómputo de Alto Desempeño}
\newacronym{acr-ood}{OOD}{fuera de distribución (\textit{Out-of-Distribution})}
\newacronym{acr-cvar}{CVaR}{valor en riesgo condicional (\textit{Conditional Value at Risk})}
\newacronym{acr-cr}{CR}{tasa de completitud de viajes (\textit{Completion Rate})}
\newacronym{acr-geh}{GEH}{estadístico empírico de comparación de flujos viales de Geoffrey E. Havers}
\newacronym{acr-iqm}{IQM}{media intercuartílica (\textit{Interquartile Mean})}
\newacronym{acr-kl}{KL}{divergencia de Kullback--Leibler}
\newacronym{acr-vf}{VF}{función de valor (\textit{Value Function})}
\newacronym{acr-clip}{CLIP}{objetivo sustituto recortado (\textit{Clipped surrogate objective})}
\newacronym{acr-posix}{POSIX}{interfaz portable de sistemas operativos (\textit{Portable Operating System Interface})}
\newacronym{acr-utm}{UTM}{sistema de coordenadas universal transversal de Mercator (\textit{Universal Transverse Mercator})}
\newacronym{acr-wgs84}{WGS84}{Sistema Geodésico Mundial 1984 (\textit{World Geodetic System 1984})}
\newacronym{acr-rht}{RHT}{sentido de circulación por la derecha (\textit{Right-Hand Traffic})}
\newacronym{acr-flir}{FLIR}{tecnología infrarroja con barrido hacia adelante (\textit{Forward Looking Infrared})}
\newacronym{acr-afac}{AFAC}{Asociación de Fábricas Argentinas de Componentes}
\newacronym{acr-fhwa}{FHWA}{Administración Federal de Carreteras (\textit{Federal Highway Administration})}
\newacronym{acr-ite}{ITE}{Instituto de Ingenieros de Transporte (\textit{Institute of Transportation Engineers})}
\newacronym{acr-conicet}{CONICET}{Consejo Nacional de Investigaciones Científicas y Técnicas}
\newacronym{acr-fcefn}{FCEFN}{Facultad de Ciencias Exactas, Físicas y Naturales}
\newacronym{acr-co}{CO}{monóxido de carbono}
\newacronym{acr-co2}{CO2}{dióxido de carbono}
\newacronym{acr-nox}{NOx}{óxidos de nitrógeno}
\newacronym{acr-hc}{HC}{hidrocarburos no combustionados}
\newacronym{acr-io}{I/O}{entrada y salida (\textit{Input/Output})}
